\section*{Hoofdstuk \ref{ch:b3:ode} --- Gewone
differentiaalvergelijkingen}

\begin{solution}{exo:b3:ode:1}
(a) Veranderlijken scheiden: $x(t) = \frac1{1 - t}$ op
$\intoo{-\infty}1$: explosie in $T_+ = 1$, met $x(t) \to
+\infty$ --- precies \cref{thm:b3:ode:maximal}(2).
(b) $x(t) = \tan t$ op $\intoo{-\pi/2}{\pi/2}$: explosie aan
beide uiteinden.
(c) $x(t) = \frac{1}{1 + \eu^{-t}}$, globaal: de oplossing blijft
in $\intoo01$, een begrensde verzameling, dus kan de grafiek niet
in eindige tijd uit elke \omterm{def:b3:topology:compact}{compacte} deelverzameling van
$\R\times\R$ ontsnappen --- $T_\pm = \pm\infty$. Merk op dat (a)
en (b) niet in tegenspraak zijn met
\cref{cor:b3:ode:globallinear}: $x^2$ en $1 + x^2$ hebben
superlineaire groei.
\end{solution}

\begin{solution}{exo:b3:ode:2}
(a) Dat is de afschatting van \cref{cor:b3:ode:uniqueness} met
$t_0 = 0$. Scherpte: voor $x' = Lx$ (globaal $L$-lipschitz)
verschillen twee oplossingen precies $(x_0 - y_0)\eu^{Lt}$.
(b) De afschatting luidt: de stromingsafbeelding op tijdstip $t$
is $\eu^{L\abs t}$-lipschitz in de beginvoorwaarde ---
\omterm{def:b3:topology:continuity}{continuïteit}, uniform voor $t$ in \omterm{def:b3:topology:compact}{compacte} verzamelingen; op
begrensde verzamelingen bij een niet globaal lipschitz $f$ voer
je hetzelfde uit op een \omterm{def:b3:topology:compact}{compacte} buis rond de \omterm{def:b3:groups:action}{banen} met de lokale
constante.
\end{solution}

\begin{solution}{exo:b3:ode:3}
(a) $\abs{\sin(tx)} \leq 1$: begrensd, dat wil zeggen lineaire
groei met $\alpha = 0$ en $\beta = 1$:
\cref{cor:b3:ode:globallinear} op $U = \R\times\R$.
(b) $\abs{tx/(1 + x^2)} \leq \abs t\cdot\frac12$: opnieuw
sublineair (zelfs begrensd op \omterm{def:b3:topology:compact}{compacte} tijdsintervallen):
globaal.
(c) $X = (x, x')$: $X' = \bigl(\begin{smallmatrix}0 & 1\\ -q(t)
& 0\end{smallmatrix}\bigr)X$: lineair met \omterm{def:b3:topology:continuity}{continue}
coëfficiënten: globaal (het kader van
\cref{thm:b3:ode:linearstructure}).
(d) Globaal; Grönwall als in \cref{cor:b3:ode:globallinear}:
$\norm{x(t)} \leq \norm{x(t_0)}\,\eu^{M\abs{t - t_0}}$ met $M =
\sup\vertiii{A}$.
\end{solution}

\begin{solution}{exo:b3:ode:4}
(a) Voor de eerste is $A^2 = -I$: $\eu^{tA} = \cos t\,I + \sin
t\,A = \bigl(\begin{smallmatrix}\cos t & -\sin t\\ \sin t &
\cos t\end{smallmatrix}\bigr)$. Jordanblok: $\lambda I$ en $N$
commuteren: $\eu^{tA} = \eu^{\lambda t}\bigl(
\begin{smallmatrix}1 & t\\ 0 & 1\end{smallmatrix}\bigr)$. Derde:
$A^2 = I$: $\eu^{tA} = \cosh t\,I + \sinh t\,A$.

(b) Stelsel $X' = \bigl(\begin{smallmatrix}0&1\\-1&0
\end{smallmatrix}\bigr)X + \bigl(\begin{smallmatrix}0\\
\cos\omega t\end{smallmatrix}\bigr)$; Duhamel met de
rotatieresolvente geeft de bijzondere oplossingen: voor $\omega
\neq 1$ is $x_p(t) = \frac{\cos(\omega t)}{1 - \omega^2}$ (ga
rechtstreeks na); voor $\omega = 1$ brengt de integraal
$\int_0^t\sin(t - s)\cos s\,\dd s = \frac t2\sin t$ de
\emph{seculiere} groei $x_p = \frac t2\sin t$ voort: resonantie
--- de aandrijving pompt energie op de eigenfrequentie en de
amplitude groeit lineair.
\end{solution}

\begin{solution}{exo:b3:ode:5}
(a) $w' = x_1x_2'' - x_1''x_2 = x_1(-px_2' - qx_2) - (-px_1' -
qx_1)x_2 = -p\,w$: dus $w(t) = w(t_0)\eu^{-\int_{t_0}^tp}$,
nooit nul of identiek nul. Is $w \neq 0$, dan zijn de vectoren
$(x_i, x_i')(t_0)$ onafhankelijk in $\R^2$, en omdat de
oplossingsruimte dimensie $2$ heeft
(\cref{thm:b3:ode:linearstructure} voor het stelsel), is $(x_1,
x_2)$ een basis.

(b) Met $u = \int\frac{\eu^{-\int p}}{x_1^2}$: $x_2 = x_1u$,
$x_2' = x_1'u + \frac{\eu^{-\int p}}{x_1}$, en
\[
x_2'' + px_2' + qx_2
= u\,(x_1'' + px_1' + qx_1) +
\Bigl(-\,p\frac{\eu^{-\int p}}{x_1} +
p\frac{\eu^{-\int p}}{x_1}\Bigr) = 0
\]
(de kruistermen vallen precies weg; werk zorgvuldig uit). Voor
$t^2x'' - 2x = 0$, dat wil zeggen $x'' - \frac2{t^2}x = 0$ ($p =
0$) met $x_1 = t^2$: $u = \int t^{-4} = -\frac1{3t^3}$, dus $x_2
= -\frac1{3t}$: algemene oplossing $at^2 + \frac bt$.
\end{solution}

\begin{solution}{exo:b3:ode:6}
Evenwichten $0$ en $1$; $f'(x) = 1 - 2x$: $f'(0) = 1 > 0$
(instabiel --- nabije oplossingen bewegen weg, zoals de
expliciete vorm toont), $f'(1) = -1 < 0$: asymptotisch \omterm{def:b3:ode:stability}{stabiel}
(\cref{thm:b3:ode:linearization} in dimensie $1$). Voor $x(0)
\in \intoo01$: daar is $f > 0$, dus stijgt de oplossing zolang ze
in $\intoo01$ blijft; ze kan nooit $0$ of $1$ bereiken
(uniciteit: dat zijn \omterm{def:b3:groups:action}{banen}), dus blijft ze daar, is ze begrensd
--- en dus globaal --- en stijgt ze naar een limiet $L \in
\intoc{x(0)}1$. Is $f(L) \neq 0$, dan is $x' \geq c > 0$ nabij de
limiet, wat $x$ voorbij $L$ dwingt: dus $f(L) = 0$ en $L = 1$;
symmetrisch gaat $x \to 0$ in $-\infty$. Expliciet bevestigt
$x(t) = \frac1{1 + C\eu^{-t}}$ alles. Algemeen: had een scalaire
autonome oplossing $x'(t_0) = 0$, dan is $x(t_0)$ een evenwicht
en maakt de uniciteit $x$ constant; anders behoudt $f(x(t))$ een
vast teken (ze verdwijnt nooit, en $t \mapsto f(x(t))$ is
\omterm{def:b3:topology:continuity}{continu}): dus is $x$ strikt monotoon --- en is een niet-constante
periodieke oplossing onmogelijk.
\end{solution}

\begin{solution}{exo:b3:ode:7}
(a) $\frac{\dd}{\dd t}H(x,y) = H_xx' + H_yy' = H_xH_y +
H_y(-H_x) = 0$.
(b) De niveauverzamelingen van $H = \frac{y^2}2 + \frac{x^4}4$
zijn \omterm{def:b3:topology:compact}{compact} ($H$ is coërcief), dus zijn de oplossingen in
\omterm{def:b3:topology:compact}{compacte} verzamelingen gevangen: globaal en begrensd
(\cref{thm:b3:ode:maximal}). Stabiliteit van $(0,0)$: $V = H$ is
positief definiet ($H = 0$ enkel in de oorsprong) met $\dot V =
0$: \cref{thm:b3:ode:lyapunov}. De linearisatie $x' = y$, $y' =
0$ heeft de niet-diagonaliseerbare nilpotente matrix met
eigenwaarde $0$: \cref{thm:b3:ode:linearization} zwijgt (haar
hypothese $\operatorname{Re}\lambda < 0$ faalt), en inderdaad is
het gelineariseerde stelsel instabiel ($y_0 \ne 0$ drijft weg)
terwijl het niet-lineaire \omterm{def:b3:ode:stability}{stabiel} is: lineariseren in een
niet-hyperbolisch evenwicht bewijst niets.
\end{solution}

\begin{solution}{exo:b3:ode:8}
(a) $\dot V = y\,y' + \sin x\cdot x' = y(-\sin x - cy) + y\sin x
= -cy^2 \leq 0$, en $V = \frac{y^2}2 + (1 - \cos x)$ is positief
definiet op $\{\abs x < 2\pi\}$ rond de oorsprong: \omterm{def:b3:ode:stability}{stabiel}
(\cref{thm:b3:ode:lyapunov}).
(b) De gelineariseerde matrix in $(0,0)$ is
$\bigl(\begin{smallmatrix}0 & 1\\ -1 &
-c\end{smallmatrix}\bigr)$, met karakteristieke veelterm
$\lambda^2 + c\lambda + 1$: wortels $\frac{-c \pm \sqrt{c^2 -
4}}2$ --- beide reëel en negatief als $c \geq 2$, complex met
reëel deel $-\frac c2 < 0$ als $0 < c < 2$. In alle gevallen is
$\operatorname{Re}\lambda < 0$:
\cref{thm:b3:ode:linearization} geeft asymptotische stabiliteit
(ondanks de ontaarde $\dot V$).
(c) In $(\pi, 0)$: $\sin(\pi + u) = -\sin u$, met linearisatie
$\bigl(\begin{smallmatrix}0&1\\1&-c\end{smallmatrix}\bigr)$ en
karakteristieke veelterm $\lambda^2 + c\lambda - 1$: wortels met
tegengesteld teken ($\lambda_+\lambda_- = -1$). Langs de
instabiele eigenvector heeft het lineaire stelsel de expliciet
ontsnappende oplossing $\eu^{\lambda_+t}v_+$ met $\lambda_+ > 0$:
de omgekeerde slinger is voor elke demping instabiel.
\end{solution}

\begin{solution}{exo:b3:ode:9}
Voor $\alpha \in \intoo01$ is naast $x \equiv 0$ elke
\[
x_c(t) = \begin{cases}0 & t \leq c,\\
\bigl((1-\alpha)(t - c)\bigr)^{1/(1-\alpha)} & t \geq c,
\end{cases}
\]
$\mathcal C^1$ en een oplossing van de vergelijking (de exponent
$\frac1{1-\alpha} > 1$ laat de afgeleide in $c$ verdwijnen): een
continuüm oplossingen door $(0,0)$. Voor $\alpha = 1$ is $x
\mapsto \abs x$ globaal $1$-lipschitz ($\abs{\abs a - \abs b}
\leq \abs{a - b}$), dus is \cref{thm:b3:ode:cauchylipschitz} van
toepassing en is de enige oplossing door $0$ $x \equiv 0$. De
grens is precies de lokale lipschitzvoorwaarde in $0$: $\abs
x^\alpha$ heeft daar onbegrensde differentiequotiënten voor
$\alpha < 1$.
\end{solution}

\begin{solution}{exo:b3:ode:10}
(a) Stel $\norm{x(t_2)} > R$ voor een $t_2 > 0$ met $\norm{x(0)}
\leq R$, en zij $t_1 = \sup\{t \leq t_2 : \norm{x(t)} \leq R\}$:
dan is $\norm{x(t_1)} = R$ en $\norm{x(t)} > R$ op
$\intoc{t_1}{t_2}$. Op dat interval heeft $g(t) =
\norm{x(t)}^2$ $g'(t) = 2\langle x(t), F(x(t))\rangle \leq 0$ (de
hypothese is van toepassing: $\norm{x(t)} \geq R$), dus $g(t_2)
\leq g(t_1) = R^2$: tegenspraak. De bal is positief invariant.
(b) Een oplossing die in de \omterm{def:b3:topology:compact}{compacte} bal gevangen zit, kan geen
$T_+ < \infty$ hebben (\cref{thm:b3:ode:maximal}(2)): globaal
voorwaarts. Gradiëntstelsel: $\frac{\dd}{\dd t}G(x(t)) =
\langle\nabla G, -\nabla G\rangle = -\norm{\nabla G(x(t))}^2
\leq 0$: $G$ daalt, dus blijft de oplossing in $\{G \leq
G(x_0)\}$, die begrensd is (coërciviteit: buiten een grote bal is
$G > G(x_0)$) en gesloten: dus \omterm{def:b3:topology:compact}{compact}. Ontsnappen is onmogelijk:
elke oplossing van een coërcief gradiëntstelsel bestaat globaal
voorwaarts en glijdt voor eeuwig bergaf.
\end{solution}

\begin{solution}{exo:b3:ode:11}
(a) Vergelijkingslemma: zij $w = x - y$ op het gemeenschappelijke
interval; dan is $w(0) \geq 0$ en $w' = x' - y' \geq F(x) - F(y)
= c(t)w$ met $c(t) = \frac{F(x) - F(y)}{x - y}$ begrensd op
\omterm{def:b3:topology:compact}{compacte} tijdsintervallen ($F$ is \omterm{def:b3:ode:lipschitz}{lokaal lipschitz}); dan is
$(w\eu^{-\int c})' \geq 0$, dus $w \geq 0$ overal. Met $F(x) =
x^2$: $y(t) = \frac1{1 - t}$ explodeert in $1$, en $x \geq y$
zolang $x$ leeft; was $T_+ > 1$, dan zou $x$ eindig zijn in $t =
1$ terwijl ze $y \to \infty$ domineert: absurd. Dus $T_+ \leq
1$.

(b) De omgekeerde vergelijking (hetzelfde lemma, met de rollen
verwisseld): op $\intcc01\cap\intco0{T_+}$ geeft $t^2 \leq 1$ dat
$x' \leq x^2 + 1$, terwijl $z(t) = \tan(t + \frac\pi4)$ voldoet
aan $z' = z^2 + 1$ met $z(0) = 1 = x(0)$: dus $x \leq z$ zolang
beide gedefinieerd zijn. Omdat $z$ eindig is op
$\intco0{\frac\pi4}$, kan $x$ niet vóór $\frac\pi4$ exploderen:
$T_+ \geq \frac\pi4$.

(c) Samen: $\frac\pi4 \leq T_+ \leq 1$ (numeriek $T_+ \approx
0.96$). Moraal: voor $x' = F(t, x)$ met $F$ superlineair in $x$
exploderen de oplossingen in eindige tijd zodra
$\int^\infty\frac{\dd s}{F(s)} < \infty$ (de vergelijkende
oplossing bereikt oneindig in die eindige tijd); lineaire groei,
waar de integraal divergeert, dwingt globaal bestaan af
(\cref{exo:b3:ode:3}). Het is dezelfde integraal van Osgood als
in \cref{pb:b3:complete:1}, nu voor het ontsnappen naar oneindig
in plaats van het ontsnappen uit nul.
\end{solution}

\begin{solution}{exo:b3:ode:12}
(a) $W' = uv'' - u''v = u(-q_2v) - (-q_1u)v = (q_1 -
q_2)\,uv$.

(b) Zij $a < b$ opeenvolgende nulpunten van $u$; normeer $u > 0$
op $\intoo ab$ (zodat $u'(a) > 0$ en $u'(b) < 0$ --- niet nul
wegens de uniciteit, want $u(a) = u'(a) = 0$ zou $u \equiv 0$
afdwingen). Stel dat $v$ geen nulpunt heeft in $\intoo ab$;
normeer $v > 0$ daar (zodat $v(a), v(b) \geq 0$ wegens de
\omterm{def:b3:topology:continuity}{continuïteit}). Integreer (a):
\[
W(b) - W(a) = \int_a^b(q_1 - q_2)\,uv \;\leq\; 0 .
\]
Maar $W(a) = u(a)v'(a) - u'(a)v(a) = -u'(a)v(a) \leq 0$ en $W(b)
= -u'(b)v(b) \geq 0$: dus $W(b) - W(a) \geq 0$. Overal
gelijkheid: $\int(q_1 - q_2)uv = 0$ met $uv > 0$ op het open
interval dwingt daar $q_1 = q_2$ af; en $W(a) = W(b) = 0$ dwingt
$v(a) = v(b) = 0$ af; dan is $W \equiv 0$ op $\intcc ab$ (haar
afgeleide verdwijnt), dat wil zeggen $(v/u)' = -W/u^2 = 0$ op
$\intoo ab$: dus $v \propto u$.

(c) Neem $q_1 = m$ en $u = \sin(\sqrt m(t - t_0))$, waarvan de
opeenvolgende nulpunten $\pi/\sqrt m$ uit elkaar liggen, en $q_2
= q \geq m$: volgens (b) verdwijnt elke oplossing $v$ van $v'' +
qv = 0$ in elk open interval van lengte $\pi/\sqrt m$ (in het
ontaarde alternatief $q \equiv m$ verdwijnt $v \propto u$
eveneens). Is daarentegen $q \leq 0$: pas (b) toe met $q_1 = q$
en $u = v$, en $q_2 = 0$ met de nulpuntvrije oplossing $\mathbf
1$ van $v'' = 0$. Had $v$ twee opeenvolgende nulpunten, dan zou
(b) ofwel een nulpunt van $\mathbf 1$ ertussen afdwingen ofwel
het ontaarde geval $\mathbf 1 \propto v$ --- beide absurd: dus
verdwijnt $v$ hoogstens één keer. Tests: voor $q = 1$ verdwijnt
$\sin t$ om de $\pi = \pi/\sqrt1$, zoals voorspeld; voor $q = -1$
verdwijnt $\sinh t$ precies één keer en $\eu^t$ nooit ---
hoogstens één nulpunt, zoals voorspeld.
\end{solution}

\begin{solution}{pb:b3:ode:1}
\textbf{1.} $\dot E = yy' + \sin x\cdot x' = -y\sin x + y\sin x =
0$: de energie is een eerste integraal. Op een \omterm{thm:b3:ode:maximal}{maximale oplossing}
is $y^2 = 2(E_0 + \cos x) \leq 2(E_0 + 1)$: $y$ is begrensd; dan
groeit $\abs{x(t)} \leq \abs{x(0)} + t\sup\abs y$ hoogstens
lineair: op elk eindig tijdsinterval blijft de \omterm{def:b3:groups:action}{baan} in een
\omterm{def:b3:topology:compact}{compacte} deelverzameling van $\R^2$, dus dwingt
\cref{thm:b3:ode:maximal}(2) $T_\pm = \pm\infty$ af. Evenwichten
$(k\pi, 0)$; linearisatie $\bigl(\begin{smallmatrix}0&1\\
\mp1&0\end{smallmatrix}\bigr)$ met $-\cos(k\pi) = \mp1$:
eigenwaarden $\pm\iu$ voor even $k$ (centrumtype, op zichzelf
onbeslissend) en $\pm1$ voor oneven $k$ (zadel).

\textbf{2.} Omdat $E$ constant is langs de oplossingen, blijft
elke \omterm{def:b3:groups:action}{baan} in een niveauverzameling $\{y^2 = 2(E_0 + \cos x)\}$.
Voor $E_0 = -1$: enkel de punten $(2k\pi, 0)$. Voor $-1 < E_0 <
1$: met $E_0 = -\cos a$ is de verzameling een disjuncte
vereniging van gesloten krommen $y = \pm\sqrt{2(\cos x - \cos
a)}$ boven $x \in [2k\pi - a, 2k\pi + a]$, één rond elk \omterm{def:b3:ode:stability}{stabiel}
evenwicht. Voor $E_0 = 1$: de krommen $y = \pm2\cos\frac x2$ die
opeenvolgende zadels verbinden --- de separatrix --- samen met de
zadels zelf. Voor $E_0 > 1$: twee grafieken $y = \pm\sqrt{2(E_0 +
\cos x)}$, voor alle $x$ gedefinieerd, die $y = 0$ nooit raken.

\textbf{3.} $V = E + 1 = \frac{y^2}2 + (1 - \cos x)$ verdwijnt in
$(0,0)$, is positief op een geperforeerde \omterm{def:b3:topology:topology}{omgeving} ($\abs x <
2\pi$), en $\dot V = 0 \leq 0$: \cref{thm:b3:ode:lyapunov} geeft
stabiliteit. Niet asymptotisch: de oplossing door $(a, 0)$ (met
$0 < a$ klein) blijft op de niveaukromme $E = -\cos a$, waarvan
de afstand tot de oorsprong positief is (de kromme ontmoet de
$x$-as enkel in $\pm a$): dus $\varphi_t(a, 0) \not\to (0,0)$.

\textbf{4.} Op de niveaukromme $C_a$ zijn er geen evenwichten
($y = 0$ dwingt $x = \pm a$ af, en $\sin(\pm a) \neq 0$ voor $0 <
a < \pi$), dus heeft de snelheid $\norm{(y, -\sin x)}$ een
positief minimum $m$ op de \omterm{def:b3:topology:compact}{compacte} $C_a$. Volg de oplossing
vanaf $(a, 0)$: in het onderste halfvlak is $x' = y < 0$, dus
daalt $x$ in eindige tijd van $a$ naar $-a$ (de integralen voor
de kwart- of halve periode convergeren: de analyse van vraag 5),
en bereikt ze $(-a, 0)$; wegens de symmetrie $(x, y) \mapsto (x,
-y)$, $t \mapsto -t$ van de vergelijking wordt de bovenste helft
in dezelfde tijd $\frac T2$ teruggelopen: de oplossing keert op
tijdstip $T$ terug naar $(a, 0)$. De uniciteit
(\cref{cor:b3:ode:uniqueness}) plant dat voort: $x(t + T) = x(t)$
voor alle $t$, dus periodiek.

\textbf{5.} Op de tak waar $y > 0$: $y = \sqrt{2(\cos x - \cos
a)}$ en $\dd t = \frac{\dd x}{y}$; $x$ van $-a$ tot $a$
integreren geeft de halve periode, en de symmetrie $x \mapsto -x$
halveert de integraal nogmaals:
\[
T(a) = 4\int_0^a\frac{\dd x}{\sqrt{2(\cos x - \cos a)}} .
\]
Convergentie in $x = a^-$: $\cos x - \cos a = \sin(a)(a - x) +
O((a-x)^2)$ met $\sin a > 0$: de integrand gedraagt zich als
$\bigl(2\sin a\,(a - x)\bigr)^{-1/2}$, \omterm{def:b3:lebesgue:l1}{integreerbaar}.

\textbf{6.} Met $\sin\frac x2 = k\sin\varphi$ en $k = \sin\frac
a2$: $\cos x - \cos a = 2(k^2 - \sin^2\frac x2) =
2k^2\cos^2\varphi$ en $\frac12\cos\frac x2\,\dd x =
k\cos\varphi\,\dd\varphi$, dus
\[
T(a) = 4\int_0^{\pi/2}
\frac{\dd\varphi}{\sqrt{1 - k^2\sin^2\varphi}} .
\]
Als $a \to 0^+$ gaat $k \to 0$: voor $k \leq k_0 < 1$ wordt de
integrand gedomineerd door $(1 - k_0^2\sin^2\varphi)^{-1/2}$,
\omterm{def:b3:topology:continuity}{continu} op $\intcc0{\pi/2}$: de gedomineerde convergentie geeft
$T \to 4\cdot\frac\pi2 = 2\pi$. Ontwikkelen van $(1 - u)^{-1/2} =
1 + \frac u2 + \frac{3u^2}8 + \cdots$ met $u = k^2\sin^2\varphi$
(normale convergentie voor $k < 1$) en $\int_0^{\pi/2}\sin^2 =
\frac\pi4$ geeft
\[
T(a) = 2\pi\Bigl(1 + \frac{k^2}{4} + O(k^4)\Bigr) :
\]
het isochronisme geldt enkel in eerste orde; de periode groeit
met de amplitude.

\textbf{7.} Als $k \uparrow 1$ stijgen de integranden naar $(1 -
\sin^2\varphi)^{-1/2} = \frac1{\cos\varphi}$, waarvan de
integraal divergeert: volgens de monotone convergentie gaat
$T(a) \to +\infty$ als $a \to \pi^-$.

\textbf{8.} Op de tak $y = 2\cos\frac x2$ ($\abs x < \pi$): $\frac
{\dd x}{2\cos(x/2)} = \dd t$; met $u = \frac x2$ is $\int\frac{\dd
u}{\cos u} = \ln\tan\bigl(\frac u2 + \frac\pi4\bigr)$, dus $t =
\ln\tan\bigl(\frac x4 + \frac\pi4\bigr)$, oftewel
\[
x(t) = 4\arctan(\eu^t) - \pi,
\qquad y(t) = x'(t) = \frac{4\eu^t}{1 + \eu^{2t}} =
\frac{2}{\cosh t} .
\]
Verificatie via de energie: met $\theta = \arctan\eu^t$ is
$\sin2\theta = \frac1{\cosh t}$, dus $\cos x = -\cos4\theta = -1
+ \frac{2}{\cosh^2t}$ en $E = \frac{y^2}2 - \cos x =
\frac2{\cosh^2t} + 1 - \frac2{\cosh^2t} = 1$: de \omterm{def:b3:groups:action}{baan} ligt op de
separatrix, en $y^2 = 2(1 + \cos x)$ differentiëren waar $y > 0$
geeft $y' = -\sin x$ terug. Limieten: $x \to \pm\pi$ en $y \to 0$
als $t \to \pm\infty$.

\textbf{9.} De \omterm{def:b3:groups:action}{baan} gaat voorwaarts naar het zadel $(\pi, 0)$ en
achterwaarts naar $(-\pi, 0)$, maar komt er nooit aan: bereikte
ze $(\pi, 0)$ op een eindig tijdstip $t^*$, dan zouden twee
verschillende \omterm{thm:b3:ode:maximal}{maximale oplossingen} --- de separatrixoplossing en
de constante oplossing in het zadel --- door hetzelfde punt
$(t^*, (\pi, 0))$ gaan, in tegenspraak met
\cref{cor:b3:ode:uniqueness}. Zadels worden enkel asymptotisch
benaderd.

\textbf{10.} Voor $E_0 > 1$: $y^2 = 2(E_0 + \cos x) \geq 2(E_0 -
1) > 0$: $y$ behoudt haar teken, en $\abs{x'} = \abs y \geq
\sqrt{2(E_0 - 1)}$: $x$ is strikt monotoon en globaal (vraag 1),
met $x(t) \to \pm\infty$. Omdat $y(t) = \pm\sqrt{2(E_0 + \cos
x(t))}$ en $\cos$ $2\pi$-periodiek is, keert $y$ telkens naar
haar waarde terug wanneer $x$ met $2\pi$ opschuift; de daarvoor
benodigde tijd is
\[
\tau(E_0) = \int_{x_0}^{x_0 + 2\pi}\frac{\dd
x}{\sqrt{2(E_0 + \cos x)}} =
\int_{-\pi}^{\pi}\frac{\dd x}{\sqrt{2(E_0 + \cos x)}}
\]
(substitutie; periodiciteit): $y$ is $\tau$-periodiek --- de
slinger wentelt rond met een voor grote energieën asymptotisch
constante rotatiesnelheid $2\pi/\tau \approx \sqrt{2E_0}$.

\textbf{11.} De methode, op volgorde: de \emph{energie} ($\dot E
= 0$) herleidt de tweedimensionale \omterm{ex:b3:ode:planeclassification}{stroming} tot eendimensionale
niveaukrommen; de \emph{begrensdheid} van $y$ op elk niveau plus
het \omterm{thm:b3:ode:maximal}{ontsnappen uit compacte verzamelingen} geeft globaal bestaan;
de \emph{\omterm{def:b3:topology:compact}{compactheid}} van de gesloten niveaus geeft
snelheidsgrenzen en dus de periodiciteit; de \emph{uniciteit}
zet de eerste terugkeer om in exacte periodiciteit, verbiedt
aankomst in een zadel in eindige tijd, en scheidt de baantypen;
\emph{linearisatie en Lyapunov} klasseren de evenwichten; en de
\emph{periode-integraal} wordt geanalyseerd met de
convergentiestellingen van \cref{ch:b3:lebesgue}. Nergens
bezaten we --- of hadden we nodig --- een gesloten algemene
oplossing: de kwalitatieve theorie haalde elk kenmerk van de
beweging uit de vergelijking zelf.

\textbf{12.} Partieel integreren: $W_n =
\int_0^{\pi/2}\sin^{2n-1}\varphi \cdot\sin\varphi\,\dd\varphi =
(2n-1)\int_0^{\pi/2} \sin^{2n-2}\varphi\cos^2\varphi\,\dd\varphi
= (2n-1)(W_{n-1} - W_n)$, dus $W_n = \frac{2n-1}{2n}W_{n-1}$; met
$W_0 = \frac\pi2$ en $\prod_{j=1}^n\frac{2j-1}{2j} =
\frac{(2n)!}{4^n(n!)^2}$ (splits $(2n)! = 2^nn!\prod(2j-1)$)
geeft inductie de weergegeven waarde. Binomiaalreeks: $(1 -
u)^{-1/2} = \sum_{n\geq0}c_nu^n$ met $c_n =
\frac{(2n)!}{4^n(n!)^2} \in \intoc01$, straal $1$. Voor $k \leq
k_0 < 1$ en $u = k^2\sin^2\varphi$ convergeert de reeks $\sum
c_nk^{2n}\sin^{2n}\varphi$ normaal in $\varphi$ (want $c_nk^{2n}
\leq k_0^{2n}$), zodat term voor term integreren in de formule
van vraag 6 geoorloofd is:
\[
T(a) = 4\sum_{n\geq0}c_nk^{2n}W_n
= 4\cdot\frac\pi2\sum_{n\geq0}c_n^2\,k^{2n}
= 2\pi\sum_{n\geq0}
\Bigl(\frac{(2n)!}{4^n(n!)^2}\Bigr)^{\!2}k^{2n} .
\]
Met $c_0 = 1$, $c_1 = \frac12$ en $c_2 = \frac38$: $T(a) =
2\pi\bigl(1 + \frac{k^2}4 + \frac{9k^4}{64} + O(k^6)\bigr)$,
waarbij de restterm uniform is voor $k \leq k_0$ (de staart wordt
door een meetkundige reeks gedomineerd).

\textbf{13.} $k = \sin\frac a2 = \frac a2 - \frac{a^3}{48} +
O(a^5)$, dus
\[
k^2 = \frac{a^2}4 - \frac{a^4}{48} + O(a^6),
\qquad
k^4 = \frac{a^4}{16} + O(a^6) ,
\]
en
\[
\frac{T(a)}{2\pi} = 1 + \frac14\Bigl(\frac{a^2}4 -
\frac{a^4}{48}\Bigr) + \frac9{64}\cdot\frac{a^4}{16} +
O(a^6)
= 1 + \frac{a^2}{16} + \frac{11\,a^4}{3072} + O(a^6) ,
\]
want $-\frac1{192} + \frac9{1024} = \frac{-16 + 27}{3072} =
\frac{11}{3072}$.

\textbf{14.} In de elliptische vorm van vraag 6 is $a \mapsto k =
\sin\frac a2$ een \omterm{def:b3:topology:continuity}{continue} strikt stijgende bijectie van
$\intoo0\pi$ op $\intoo01$, en voor elke $\varphi \in
\intoc0{\pi/2}$ is de integrand $(1 - k^2\sin^2\varphi)^{-1/2}$
strikt stijgend in $k$: dus is $T$ strikt stijgend.
\omterm{def:b3:topology:continuity}{Continuïteit}: op $k \leq k_0 < 1$ wordt de integrand gedomineerd
door de \omterm{def:b3:topology:continuity}{continue} $(1 - k_0^2\sin^2\varphi)^{-1/2}$, zodat de
gedomineerde convergentie langs $k' \to k$ van toepassing is. Met
de limieten $T \to 2\pi$ als $a \to 0^+$ (vraag 6) en $T \to
+\infty$ als $a \to \pi^-$ (vraag 7) maken de strikte monotonie
en de tussenwaardestelling van $T$ een bijectie van $\intoo0\pi$
op $\intoo{2\pi}{+\infty}$.

\textbf{15.} Een klok telt slingerbewegingen; geregeld bij
verdwijnende amplitude boekt ze de harmonische periode $2\pi$ per
beweging (in de tijdseenheid van de slinger). Bij amplitude $a$
is de ware periode $T(a) = 2\pi\bigl(1 + \frac{a^2}{16} +
O(a^4)\bigr)$: de klok boekt $2\pi$ terwijl er werkelijk $T(a)$
verstrijkt, dus loopt ze achter met de fractie
\[
\frac{T(a) - 2\pi}{T(a)} = \frac{a^2}{16} + O(a^4) .
\]
Voor $a = 0.2$ rad (ongeveer $11.5$ graden): $\frac{a^2}{16} =
\frac{0.04}{16} = \frac1{400}$, en een dag telt $86400$ s: de
klok verliest $86400/400 = 216$ seconden --- een goede drieënhalve
minuut --- per dag. Vandaar de twee historische remedies: een
minieme constante amplitude afdwingen (het echappement), of de
dwangvoorwaarde zo buigen dat de periode exact
amplitude-onafhankelijk wordt (de cycloïdale wangen van Huygens,
1657).

\textbf{16.} In $\tau(E_0) = \int_{-\pi}^{\pi}\frac{\dd
x}{\sqrt{2(E_0 + \cos x)}}$ is de integrand voor elke vaste $x$
strikt dalend in $E_0$: dus is $\tau$ strikt dalend. Als $E_0
\downarrow 1$ stijgen de integranden puntsgewijs naar $\bigl(2(1
+ \cos x)\bigr)^{-1/2} = \frac1{2\abs{\cos\frac x2}}$, waarvan de
integraal over $\intoo{-\pi}\pi$ divergeert ($\cos\frac x2$
verdwijnt in $\pm\pi$ in eerste orde): de monotone convergentie
geeft $\tau(E_0) \to +\infty$. Als $E_0 \to +\infty$:
\[
\sqrt{2E_0}\,\tau(E_0) = \int_{-\pi}^{\pi}
\frac{\dd x}{\sqrt{1 + \cos x/E_0}} \longrightarrow 2\pi
\]
volgens de gedomineerde convergentie (voor $E_0 \geq 2$ is de
integrand hoogstens $\sqrt2$). Dus $\tau \approx
2\pi/\sqrt{2E_0}$: één omwenteling duurt zo lang als vrije
rotatie met snelheid $\sqrt{2E_0}$, met de potentiaal gereduceerd
tot een rimpeling --- in overeenstemming met de rotatiesnelheid
uit vraag 10.

\textbf{17.} De assen dragen de expliciete oplossingen $t \mapsto
(x_0\eu^{\alpha t}, 0)$ en $t \mapsto (0, y_0\eu^{-\gamma t})$,
samen met het evenwicht $(0, 0)$: ze zijn verenigingen van \omterm{def:b3:groups:action}{banen}.
Het veld is $\mathcal C^1$, dus \omterm{def:b3:ode:lipschitz}{lokaal lipschitz}; een oplossing
die in $Q$ start en een as zou raken, zou door een punt van een
van die \omterm{def:b3:groups:action}{banen} gaan en er volgens
\cref{cor:b3:ode:uniqueness} mee samenvallen --- onmogelijk, want
de ene leeft op de as en de andere niet. Dus is $Q$ invariant in
beide tijdsrichtingen. Evenwichten in $Q$: $x > 0$ dwingt $\alpha
- \beta y = 0$ af en $y > 0$ dwingt $\delta x - \gamma = 0$ af:
het enkele punt $(x_*, y_*) = \bigl(\frac\gamma\delta,
\frac\alpha\beta\bigr)$.

\textbf{18.} Langs een oplossing is
\[
\dot H = \Bigl(\delta - \frac\gamma x\Bigr)x' +
\Bigl(\beta - \frac\alpha y\Bigr)y'
= (\delta x - \gamma)(\alpha - \beta y) +
(\beta y - \alpha)(\delta x - \gamma) = 0 .
\]
$f(x) = \delta x - \gamma\ln x$ heeft $f'' = \gamma/x^2 > 0$, een
$f'$ die enkel in $x_*$ verdwijnt, en $f \to +\infty$ zowel in
$0^+$ als in $+\infty$: strikt convex en eigenlijk, met minimum
$f(x_*)$; en evenzo $g(y) = \beta y - \alpha\ln y$, met minimum
$g(y_*)$. Dus is $H \geq h_*$ met gelijkheid enkel in $(x_*,
y_*)$, en is elke sublevelverzameling $\{H \leq h\}\cap Q$
\omterm{def:b3:topology:compact}{compact}: $f(x) \leq h - g(y_*)$ sluit $x$ wegens de
eigenlijkheid op in een \omterm{def:b3:topology:compact}{compact} interval van $\intoo0{+\infty}$,
en evenzo $y$, en de verzameling is gesloten in $\R^2$ omdat $H
\to +\infty$ aan de rand van $Q$. Een \omterm{thm:b3:ode:maximal}{maximale oplossing} blijft
op haar \omterm{def:b3:topology:compact}{compacte} niveauverzameling, dus kan ze niet in eindige
tijd elke \omterm{def:b3:topology:compact}{compacte} verzameling verlaten:
\cref{thm:b3:ode:maximal} maakt haar globaal.

\textbf{19.} Leg $h > h_*$ vast en stel $c = h - g(y_*) >
f(x_*)$. Omdat $f$ strikt daalt van $+\infty$ tot $f(x_*)$ op
$\intoc0{x_*}$ en strikt terugstijgt tot $+\infty$ op
$\intco{x_*}{+\infty}$, heeft de vergelijking $f(x) = c$ precies
twee wortels $x_- < x_* < x_+$, en is $\{f \leq c\} =
\intcc{x_-}{x_+}$. Voor $x \in \intoo{x_-}{x_+}$ heeft $g(y) = h
- f(x) > g(y_*)$ precies twee wortels $y_-(x) < y_* < y_+(x)$,
\omterm{def:b3:topology:continuity}{continu} in $x$ (inversen van de strikt monotone \omterm{def:b3:topology:continuity}{continue}
beperkingen van $g$ aan weerszijden van $y_*$), met $y_\pm(x) \to
y_*$ als $x \to x_\pm$; in $x = x_\pm$ is de enige oplossing $y =
y_*$. Dus is $\{H = h\}\cap Q$ de vereniging van de grafieken van
$y_+$ en $y_-$ boven $\intcc{x_-}{x_+}$, aan elkaar gelijmd in
$(x_\pm, y_*)$: een gesloten kromme rond $(x_*, y_*)$ --- het
analogon van de ovalen van de slinger.

\textbf{20.} Zij $C_h = \{H = h\}\cap Q$ met $h > h_*$: het enige
evenwicht van $Q$ ligt niet op $C_h$, dus verdwijnt het veld daar
nergens. Tekens: $x' = \beta x(y_* - y)$, $y' = \delta y(x -
x_*)$: de beweging gaat naar rechts onder de rechte $y = y_*$,
omhoog rechts van $x = x_*$, naar links erboven en omlaag links
--- circulatie tegen de wijzers van de klok in. Volg de oplossing
vanaf een punt $(x_0, y_-(x_0))$ van de open onderste tak, waar
$x' > 0$: de tijd om de rechterhoek $B = (x_+, y_*)$ te bereiken
is
\[
\int_{x_0}^{x_+}\frac{\dd x}{\beta x\,\bigl(y_* -
y_-(x)\bigr)} .
\]
Kies nabij $x_+$ een $x' \in \intoo{x_*}{x_+}$; voor $x \in
\intcc{x'}{x_+}$ is $f(x_+) - f(x) \geq f'(x')\,(x_+ - x)$ ($f'$
is stijgend en positief voorbij $x_*$), terwijl de
niveaubetrekking en de ongelijkheid van Taylor $g(y_-(x)) -
g(y_*) \leq \frac12\,\bigl(\max g''\bigr)\,(y_* - y_-(x))^2$
geven op het \omterm{def:b3:topology:compact}{compacte} $y$-bereik van $C_h$: dus $y_* - y_-(x)
\geq c\,\sqrt{x_+ - x}$ met $c > 0$, en de integrand is
$O\bigl((x_+ - x)^{-1/2}\bigr)$: \omterm{def:b3:lebesgue:l1}{integreerbaar} --- de convergentie
van vraag 5, overgezet. Elders op de tak is de integrand \omterm{def:b3:topology:continuity}{continu}.
Dus wordt $B$ in eindige tijd bereikt; daar is $y' = \delta
y_*(x_+ - x_*) > 0$, gaat de \omterm{def:b3:groups:action}{baan} het gebied $x > x_*$, $y > y_*$
binnen, klimt ze met de symmetrische afschatting (de rollen van
$f$ en $g$ verwisseld) naar de bovenhoek $(x_*, y_+)$, enzovoort
rond de vier bogen: na een eindige tijd $T > 0$ keert de
oplossing terug naar haar beginpunt. Volgens
\cref{cor:b3:ode:uniqueness} is ze $T$-periodiek --- het argument
van vraag 4, woordelijk.

\textbf{21.} Op een $T$-periodieke \omterm{def:b3:groups:action}{baan} in $Q$ is $t \mapsto \ln
x(t)$ $\mathcal C^1$ en $T$-periodiek, dus
\[
0 = \int_0^T(\ln x)'\,\dd t
= \int_0^T(\alpha - \beta y)\,\dd t
= \alpha T - \beta\int_0^Ty\,\dd t ,
\]
wat $\frac1T\int_0^Ty = \frac\alpha\beta$ geeft; en evenzo geeft
$0 = \int_0^T(\ln y)' = \delta\int_0^Tx\,\dd t - \gamma T$ dat
$\frac1T\int_0^Tx = \frac\gamma\delta$. De tijdsgemiddelden zijn
de evenwichtswaarden, voor elke \omterm{def:b3:groups:action}{baan} ongeacht de amplitude ---
een behoudswet die niemand met de hand heeft ingelegd.

\textbf{22.} Met oogst is het stelsel opnieuw van
lotka-volterravorm, met parameters $\alpha - \varepsilon$,
$\beta$, $\gamma + \varepsilon$, $\delta$ (het inwendige
evenwicht blijft bestaan omdat $\varepsilon < \alpha$). Vraag 21
toegepast op het nieuwe stelsel:
\[
\bar x = \frac{\gamma + \varepsilon}{\delta}
\ \ (\text{het prooigemiddelde stijgt}),
\qquad
\bar y = \frac{\alpha - \varepsilon}{\beta}
\ \ (\text{het roofdiergemiddelde daalt}) :
\]
ongericht oogsten verschuift het evenwicht naar de prooi. De
gegevens van d'Ancona lezen dat achterstevoren: de oorlog sneed
de visserij terug, $\varepsilon$ daalde, dus steeg het
roofdiergemiddelde $(\alpha - \varepsilon)/\beta$ en daalde het
prooigemiddelde --- een groter haaienaandeel in de vangst,
precies wat de Adriatische vismarkten optekenden. Dat is het
beginsel van Volterra, hetzelfde mechanisme achter de
pesticideparadoxen: beide trofische niveaus uitdunnen komt het
niveau ten goede dat wordt opgegeten.

\textbf{23.} Het recept, beide keren: (i) een eerste integraal
--- $E$ voor de slinger, hier $H$, gevonden door $\dd y/\dd x$ te
scheiden --- klapt het vlak samen tot krommen; (ii) de
eigenlijkheid en de \omterm{def:b3:topology:compact}{compactheid} van de niveauverzamelingen geven
globaal bestaan via \cref{thm:b3:ode:maximal}; (iii) de
meetkunde van de niveaus --- ovalen, daar uit de vorm van $\cos$
en hier uit de strikte convexiteit van $f$ en $g$ --- wordt van
de integraal afgelezen, niet van de \omterm{ex:b3:ode:planeclassification}{stroming}; (iv) een nergens
verdwijnend veld op een \omterm{def:b3:topology:compact}{compacte} ovaal plus \omterm{def:b3:lebesgue:l1}{integreerbare}
singulariteiten in de hoeken dwingt een eindige terugkeertijd af;
(v) de uniciteit (\cref{cor:b3:ode:uniqueness}) zet de terugkeer
om in periodiciteit en verbiedt aankomst in een evenwicht in
eindige tijd; (vi) de dividenden --- periodeontwikkelingen,
wetten van gemiddelden --- komen van de convergentiestellingen
toegepast op de resulterende integralen. Noch $x'' = -\sin x$
noch Lotka--Volterra laat een elementaire gesloten oplossing toe
(elliptische integralen in het ene geval, transcendente
niveaukrommen in het andere), en nergens was er een nodig: de
vergelijking zelf, kwalitatief ondervraagd, gaf de hele beweging
prijs.

\textbf{24.} Op een energieniveau $E_0 > 1$ is $y^2 = 2(E_0 +
\cos x) > 0$: de extremen van $y^2$ liggen bij $\cos x = \pm1$,
wat de vermelde grenzen geeft, bereikt in $x \equiv 0, \pi$.
Tijdsgemiddelde over één periode $\tau = \tau(E_0)$: $x$ schuift
precies $2\pi$ op, dus
\[
\frac1\tau\int_0^\tau y\,\dd t = \frac{x(\tau) -
x(0)}{\tau} = \frac{2\pi}\tau .
\]
De verhouding van de extreme snelheden is $\sqrt{\frac{E_0 +
1}{E_0 - 1}} = 1 + O(E_0^{-1}) \to 1$: bij hoge energie is de
rimpeling $\pm1$ van de potentiaal verwaarloosbaar tegenover
$E_0$, en draait de slinger bijna uniform rond --- het wasbord
vlakt af.

\textbf{25.} Op $E_0 \geq 1 + \delta$ wordt de integrand van
$\tau(E_0) = \int_{-\pi}^{\pi}\frac{\dd x}{\sqrt{2(E_0 + \cos
x)}}$ gedomineerd door $(2\delta)^{-1/2}$ en haar afgeleide naar
$E_0$
\[
\partial_{E_0}\frac1{\sqrt{2(E_0 + \cos x)}}
= -\frac{1}{\bigl(2(E_0 + \cos x)\bigr)^{3/2}}
\]
door $(2\delta)^{-3/2}$, beide \omterm{def:b3:lebesgue:l1}{integreerbaar} op de \omterm{def:b3:topology:compact}{compacte}
$\intcc{-\pi}\pi$: differentiëren onder de integraal
(\cref{thm:b3:lebesgue:paramdiff}) is dus van toepassing en geeft
$\tau'(E_0) < 0$ (de integrand is strikt negatief): strikt
dalend en $\mathcal C^1$. Limieten: als $E_0 \to \infty$ is $\tau
\leq \frac{2\pi}{\sqrt{2(E_0 - 1)}} \to 0$; als $E_0 \to 1^+$
schrijf je $E_0 + \cos x = (E_0 - 1) + 2\cos^2\frac x2$; de
integrand stijgt als $E_0$ daalt, dus geeft de monotone
convergentie
\[
\tau(E_0) \;\nearrow\;
\int_{-\pi}^{\pi}\frac{\dd x}{2\,\abs{\cos\frac x2}}
= +\infty,
\]
waarbij de limietintegraal in $x = \pm\pi$ divergeert (daar is
$\abs{\cos\frac x2} \sim \frac{\abs{x \mp \pi}}2$, een niet
\omterm{def:b3:lebesgue:l1}{integreerbare} $\frac1{\abs\cdot}$): de periode explodeert bij het
naderen van de separatrix, in overeenstemming met $T(a) \to
\infty$ van de libratiekant uit Deel II. De energieas leest zich
zo: rust bij $E_0 = -1$; libraties, $2\pi \nearrow \infty$, op
$\intoo{-1}1$; de oneindig trage separatrix bij $E_0 = 1$; en
daarbuiten rotaties, $\infty \searrow 0$. Eén integraal, het hele
leven van de slinger.
\end{solution}
